\section{Conclusion}
\label{sec:conclusion}

This volume has shown that the structure of the Standard Model
can be interpreted as arising from closure geometry.
Beginning from the closure energy functional and the minimal condition
for structural persistence, we have identified:

\begin{itemize}
  \item the triangle as the unique minimal closure unit
        (Section~\ref{sec:triangular});
  \item $\mathrm{SU}(3)\times\mathrm{SU}(2)\times\mathrm{U}(1)$
        as the symmetry structure of closure invariance
        (Section~\ref{sec:symmetry});
  \item particle classes as structural regimes of triangular degeneracy
        (Section~\ref{sec:particles});
  \item generations as sequential stabilization of triangular edges
        (Section~\ref{sec:generations});
  \item electric charge as phase asymmetry within closure
        (Section~\ref{sec:charge});
  \item the weak interaction as internal fluctuation of degenerate closure
        (Section~\ref{sec:weak});
  \item mass and the Higgs mechanism as structural fixing of asymmetry
        (Section~\ref{sec:higgs}).
\end{itemize}

The overall logical chain is:
\begin{equation}
  \text{difference}
  \;\to\; C_{ij}
  \;\to\; \text{triangular closure}
  \;\to\; \mathrm{SU}(3)\times\mathrm{SU}(2)\times\mathrm{U}(1)
  \;\to\; \text{Standard Model.}
\end{equation}

The Standard Model can be interpreted as a structural consequence
of closure geometry.
This interpretation is qualitative in nature
and does not yet constitute a predictive theory.
The path toward quantitative consequences and large-scale geometry
is taken up in Vol.~3.
